\documentclass[10pt,a4paper]{amsart}
\usepackage[margin=19mm]{geometry}
\usepackage{amsmath,amssymb,mathtools}
\usepackage[scr=rsfs,cal=euler]{mathalfa}
\usepackage{xfrac}
\usepackage[unicode,hidelinks,bookmarks=false]{hyperref}
\newcommand{\ie}{i.e.\ }
\newcommand{\cf}{cf.\ }
\newcommand{\resp}{respectively\ }
\newcommand{\Subsection}{\subsection}
\providecommand{\nobreakdash}{\nobreak\hbox{-}}
\setlength{\emergencystretch}{2em}
\multlinegap0pt
\usepackage{bm}
\usepackage{bbm}
\usepackage{dsfont}
\usepackage{xspace}
\usepackage{cancel}
\usepackage{mathtools}
\usepackage{amsthm}
\makeatletter
\@ifundefined{prop}{\newtheorem{prop}{Proposition}}{}
\makeatother
\newcommand{\wagrad}{\nabla \frac{\delta J}{\delta \rho}}
\title[Implicit Bias of the JKO Scheme]{Implicit Bias of the JKO Scheme}
\author{Peter Halmos, Boris Hanin}
\date{Source: arXiv:2511.14827v3 (CC BY 4.0)}
\begin{document}
\maketitle

\noindent\emph{Source and licence.} Implicit Bias of the JKO Scheme. Version arXiv:2511.14827v3 (CC BY 4.0), distributed under the Creative Commons Attribution 4.0 International licence, as recorded in the arXiv licence metadata for this exact version and shown on the abstract page https://arxiv.org/abs/2511.14827v3. Full rights notice: licenses/arxiv-2511.14827-CC-BY-4.0.txt.\par\smallskip

\noindent\textit{Editorial scope.} Counted unit: the Prop whose source label is \texttt{None} in the article (source0.tex lines 2188--2201, source label \texttt{None}), delivered together with the article's own complete original proof of it (source0.tex lines 2202--2255, one \texttt{proof} environment of 54 lines). No mathematics is written, repaired or re-derived: statement and proof are the article's own text line for line. Required context retained uncounted: prop:grad\_first\_var (lines 2138--2143). Every definition, display and label of those lines is retained unchanged. Omitted from this package and not redistributed: 1--2137, 2144--2187, 2256--2671. Nothing inside a retained excerpt is omitted or altered: every retained body is delivered exactly as the article's lines are written, so no omission receipt of any kind accompanies this package. Citations: the retained excerpts carry no citation command at all, so no bibliography entry of the article is transcribed and no reference is redistributed. Reference adaptation: none was needed and none was made; every reference inside the delivered unit resolves inside it, so no unresolved reference or citation is produced. Printed slips kept literal: no printed word, symbol, hypothesis or number of the counted statement or of its proof was altered. Layout and class: the article's own class is \texttt{amsart} with packages including inputenc, fontenc, geometry, amsmath, amssymb, amsfonts, amsthm, mathtools, euscript, latexsym, bm, graphicx, xcolor, tikz; the wrapper is the lane's pinned \texttt{amsart} editorial wrapper at 10pt on A4, which restates the theorem-like environments the retained text needs and adds the article's own macro definitions that the retained text needs (\texttt{\textbackslash wagrad}), with extra wrapper packages (bm, bbm, dsfont, xspace, cancel, mathtools). The delivered numbering is the unit's own. Assets: no retained span contains a figure environment, a graphics-inclusion command, a PDF-page inclusion, a table or a TikZ picture; no image file is copied into this package, the package declares no asset dependency and no image file is redistributed with it. Licence: version arXiv:2511.14827v3 of this article is distributed under the Creative Commons Attribution 4.0 International licence, as recorded in the arXiv licence metadata for this exact version and shown on the abstract page https://arxiv.org/abs/2511.14827v3. A full rights notice is delivered at licenses/arxiv-2511.14827-CC-BY-4.0.txt. Characters: every retained excerpt is pure ASCII, so the delivered engine, XeLaTeX with its default Latin Modern text font, reports no missing character.
\begin{prop}\label{prop:grad_first_var}
Let $\rho \in {P}_{ac}(\mathbb{R}^{d})$ be a solution to the continuity equation $\partial_{t} \rho =- \nabla \cdot (\rho(x,t) \mathbf{v}(x,t))$ for a given velocity field $\mathbf{v}_{t}(x): [0,1] \times \mathbb{R}^{d}  \to \mathbb{R}^{d}$ in $T_{\rho} P_{ac}(\mathbb{R}^{d})$ (i.e. a gradient field) and a fixed measure $\rho^{(n,\eta)}$. Then;
\begin{align*}
    \nabla \frac{\delta}{\delta \rho} \int \frac{\delta W_{2}^{2}(\rho^{(n,\eta)}, \rho) }{\delta \rho} -\nabla\cdot \left( \rho_{t} \mathbf{v}_t \right) = 2 \,\mathbf{v}_t
\end{align*}
\end{prop}

\begin{prop}
    Let $\rho^{(n,\eta)} \in {P}_{ac}(\mathbb{R}^{d})$ be a fixed measure, and suppose $\mathbf{v}_{t}: [0,1] \times \mathbb{R}^{d}  \to \mathbb{R}^{d}$ is a time-varying vector field. Assume the following boundary conditions holds at the domain boundaries
    \[ \left. \left( \partial_{t} ( \rho_{t} \mathbf{v}_{t} ) \cdot \bm{n}\right) \right|_{\partial \Omega}  = 0, \quad \left. \rho_{t} \right|_{\partial \Omega} = 0\]
so that not only does the density $\rho$ vanish on the boundary of the domain, but also the normal component $\bm{n}$ of the time-change of the mass-flux $\rho \mathbf{v}_{t}$. Then, one has the following integral over the partial in time $\partial_{t}$ of the mass-flux $(\rho_{t} \mathbf{v}_{t}  
    )$;
    \[
    \frac{1}{2}\nabla\frac{\delta}{\delta\rho} \int \frac{\delta W_{2}^{2}(\rho^{(n,\eta)}, \rho) }{\delta \rho} - \nabla\cdot\left( \partial_{t} (\rho_{t} \mathbf{v}_{t}  
    ) \right) = \partial_{t} \mathbf{v}_{t} +\frac{1}{2}\nabla \left\lVert \mathbf{v}_{t} \right\rVert_{2}^{2}
    \]
    and if $\mathbf{v}_{t} = - \wagrad (\rho_{t})(x(t)) := - \wagrad (t, x(t))$ is the Wasserstein gradient flow velocity, then
    \[
    \frac{1}{2}\nabla\frac{\delta}{\delta\rho} \int \frac{\delta W_{2}^{2}(\rho^{(n,\eta)}, \rho) }{\delta \rho} - \nabla\cdot\left( \partial_{t}  \left( - \wagrad (t, x(t)) \rho_{t}\right) \right) = - \partial_{t} \wagrad + \frac{1}{2}\nabla \left\lVert \wagrad \right\rVert_{2}^{2}
    \]
\end{prop}\label{prop:accn_HJ}

\begin{proof}
First, observe that
\begin{align*}
\nabla\frac{\delta}{\delta\rho} \int \frac{\delta W_{2}^{2}(\rho^{(n,\eta)}, \rho) }{\delta \rho} - \nabla\cdot\left( \partial_{t} (\rho_{t} \mathbf{v}_{t}  
    ) \right) = \nabla\frac{\delta}{\delta\rho} \int \frac{\delta W_{2}^{2}(\rho^{(n,\eta)}, \rho) }{\delta \rho} - \nabla\cdot\left( \partial_{t} \rho_{t} \mathbf{v}_{t}  + \rho_{t} \partial_{t} \mathbf{v}_{t}
     \right)
\end{align*}
By linearity of the divergence this equals the two terms
\begin{align*}
    = \nabla\frac{\delta}{\delta\rho} \int \frac{\delta W_{2}^{2}(\rho^{(n,\eta)}, \rho) }{\delta \rho} - \nabla\cdot\left( \rho_{t} \partial_{t} \mathbf{v}_{t}
     \right) + \nabla\frac{\delta}{\delta\rho} \int \frac{\delta W_{2}^{2}(\rho^{(n,\eta)}, \rho) }{\delta \rho} - \nabla\cdot\left( \partial_{t} \rho_{t} \mathbf{v}_{t} 
     \right)
\end{align*}
by Proposition~\ref{prop:grad_first_var}, we have that since $- \nabla\cdot\left( \rho_{t} \partial_{t} \mathbf{v}_{t}
     \right)$ is simply a continuity equation driven by the velocity $\partial_{t} \mathbf{v}_{t}$, the first term immediately becomes
\begin{align*}
    \nabla\frac{\delta}{\delta\rho} \int \frac{\delta W_{2}^{2}(\rho^{(n,\eta)}, \rho) }{\delta \rho} - \nabla\cdot\left( \rho_{t} \partial_{t} \mathbf{v}_{t}
     \right) = 2 \partial_{t} \mathbf{v}_{t}
\end{align*}
The other term we analyze independently, as it does not reduce to a simple integral against a continuity equation. The boundary conditions on the flux directly imply that we may apply integration by parts on terms of the form $\partial_{tt} \rho_{t}  = - \partial_{t}( \nabla \cdot (\rho_{t} \mathbf{v}_{t} ))  = -\nabla \cdot (\partial_{t}(\rho_{t} \mathbf{v}_{t}))$ and thus $-\nabla \cdot (-\nabla \cdot\left(\rho_{t} \mathbf{v}_{t}\right) \mathbf{v}_{t} ) $ as the condition
\[
\partial_{t} (\rho_{t} \mathbf{v}_{t}) \cdot \bm{n} = \underbrace{-\nabla \cdot\left(\rho_{t} \mathbf{v}_{t}\right) \mathbf{v}_{t} \cdot \bm{n} + \rho_{t} \partial_{t}\mathbf{v}_{t} \cdot \bm{n} = -\nabla \cdot\left(\rho_{t} \mathbf{v}_{t}\right) \mathbf{v}_{t} \cdot \bm{n} +0 = 0}_{\text{on }\partial \Omega}
\]
holds at the boundary $\partial \Omega$. Thus, returning to the second term
\begin{align*}
    & \nabla\frac{\delta}{\delta\rho} \int \frac{\delta W_{2}^{2}(\rho^{(n,\eta)}, \rho) }{\delta \rho} - \nabla\cdot\left( \partial_{t} \rho_{t} \mathbf{v}_{t}) \right) 
\end{align*}
we apply integration by parts to find
\begin{align*}
    &= +\nabla\frac{\delta}{\delta\rho} \int \left\langle \nabla\frac{\delta W_{2}^{2}(\rho^{(n,\eta)}, \rho) }{\delta \rho}, \mathbf{v}_{t} \right\rangle \partial_{t} \rho_{t} \\
    &= -\nabla\frac{\delta}{\delta\rho} \int \left\langle \nabla\frac{\delta W_{2}^{2}(\rho^{(n,\eta)}, \rho) }{\delta \rho}, \mathbf{v}_{t} \right\rangle \nabla \cdot (\rho_{t} \mathbf{v}_{t} )
    \\
    &= + \nabla\frac{\delta}{\delta\rho} \int \left\langle \mathbf{v}_{t}, \nabla\left\langle \nabla\frac{\delta W_{2}^{2}(\rho^{(n,\eta)}, \rho) }{\delta \rho}, \mathbf{v}_{t} \right\rangle \right\rangle d\rho_{t} \\
    &= \nabla\frac{\delta}{\delta\rho} \int \left\langle \mathbf{v}_{t}, \nabla^{2}\frac{\delta W_{2}^{2}(\rho^{(n,\eta)}, \rho) }{\delta \rho} \mathbf{v}_{t} + \nabla \mathbf{v}_{t} \nabla \frac{\delta W_{2}^{2}(\rho^{(n,\eta)}, \rho) }{\delta \rho} \right\rangle  d\rho_{t} \\
    &= \nabla\frac{\delta}{\delta\rho} \int \left\langle \mathbf{v}_{t}, \nabla^{2}\frac{\delta W_{2}^{2}(\rho^{(n,\eta)}, \rho) }{\delta \rho} \mathbf{v}_{t}  \right\rangle  d\rho_{t} + \nabla\frac{\delta}{\delta\rho} \int \left\langle \mathbf{v}_{t},  \nabla \mathbf{v}_{t} \nabla \frac{\delta W_{2}^{2}(\rho^{(n,\eta)}, \rho) }{\delta \rho} \right\rangle  d\rho_{t} 
    \end{align*}
    From Proposition~\ref{prop:grad_first_var}, we have that this simply becomes
    \begin{align*}
        &\nabla\frac{\delta}{\delta\rho} \int \left\langle \mathbf{v}_{t}, \nabla\mathbf{v}_{t} \mathbf{v}_{t}  \right\rangle  d\rho_{t} + \nabla\frac{\delta}{\delta\rho} \int \left\langle \mathbf{v}_{t},  \nabla \mathbf{v}_{t} \mathbf{v}_{t} \right\rangle  d\rho_{t} = 2 \nabla\frac{\delta}{\delta\rho} \int \left\langle \mathbf{v}_{t}, \frac{1}{2}\nabla \lVert \mathbf{v}_{t} \rVert_{2}^{2} \right\rangle  d\rho_{t} \\
        &=  \nabla\frac{\delta}{\delta\rho} \int \left\langle \mathbf{v}_{t}, \nabla \lVert \mathbf{v}_{t} \rVert_{2}^{2} \right\rangle  d\rho_{t}  =\nabla\frac{\delta}{\delta\rho} \int \lVert \mathbf{v}_{t} \rVert_{2}^{2}  d \left(- \nabla \cdot (\rho_{t} \mathbf{v}_{t} ) \right) = \nabla\frac{\delta}{\delta\rho} \int \lVert \mathbf{v}_{t} \rVert_{2}^{2}  d (\partial_{t} \rho_{t}) \\
        &= \nabla \lVert \mathbf{v}_{t} \rVert_{2}^{2}
    \end{align*}
    Thus, we have that 
    \begin{align*}
        \nabla\frac{\delta}{\delta\rho} \int \frac{\delta W_{2}^{2}(\rho^{(n,\eta)}, \rho) }{\delta \rho} - \nabla\cdot\left( \partial_{t} (\rho_{t} \mathbf{v}_{t}  
    ) \right) = 2 \left( \partial_{t} \mathbf{v}_{t} +\frac{1}{2}\nabla \left\lVert \mathbf{v}_{t} \right\rVert_{2}^{2} \right)
    \end{align*}
    And for $\mathbf{v}_{t} = - \wagrad(\rho_{t}) (x) := - \wagrad(t, x)$ this becomes
    \begin{align*}
        \nabla\frac{\delta}{\delta\rho} \int \frac{\delta W_{2}^{2}(\rho^{(n,\eta)}, \rho) }{\delta \rho} - \nabla\cdot\left( \partial_{t} (\rho_{t} -\wagrad  
    ) \right) = 2 \left( - \partial_{t} \wagrad + \frac{1}{2}\nabla \left\lVert \wagrad \right\rVert_{2}^{2} \right)
    \end{align*}
    So that we conclude.
\end{proof}

\end{document}
